\documentclass[11pt,a4paper]{article}

\usepackage[utf8]{inputenc}
\usepackage[T1]{fontenc}
\usepackage{lmodern}
\usepackage{amsmath,amssymb}
\usepackage{booktabs}
\usepackage{graphicx}
\usepackage{hyperref}
\usepackage{geometry}
\usepackage{xcolor}

\geometry{margin=2.5cm}
\hypersetup{colorlinks=true,linkcolor=blue!50!black,citecolor=blue!50!black,urlcolor=blue!50!black}

\newcommand{\Rset}{\mathcal{R}}
\newcommand{\Dset}{\mathcal{D}}
\newcommand{\wvec}{\mathbf{w}}
\newcommand{\xvec}{\mathbf{x}}

\title{\textbf{Ordinal Logistic Regression with Split Conformal Prediction}\\[0.4em]
\large for Clothing Fit Recommendation\\
\large A classical machine learning pipeline with calibrated uncertainty}

\author{Patrick Oliveira}
\date{\today}

\begin{document}

\maketitle

\begin{abstract}
We predict clothing fit (\emph{small}, \emph{fit}, \emph{large}) as an
\emph{ordinal} classification problem using a proportional-odds cumulative
link model, trained by maximizing the likelihood with
\texttt{scipy.optimize} -- classical machine learning, with no neural
network anywhere in the pipeline. We wrap the fitted model in split
conformal prediction, implemented from scratch in NumPy, to turn a bare
point prediction into a calibrated \emph{set} of plausible fit outcomes with
a distribution-free marginal coverage guarantee. Users are split at the
group level (train/validation/test share no user) so the model is evaluated
on genuinely unseen shoppers, not on memorized body measurements. On a
synthetic fixture matching the real dataset's clustering and class skew, the
ordinal model attains a mean absolute ordinal error of
$0.3535$ versus $0.4397$ for a majority-class floor and
$0.6298$ for an unordered multinomial logistic regression baseline with
identical features, and the conformal wrapper achieves
$88.55\%$ empirical coverage against a $90\%$ target. We state explicitly
why standard ranking metrics (NDCG, Recall@K, MAP) do not apply to this
problem.
\end{abstract}

\tableofcontents
\newpage

\section{Notation}
\label{sec:notation}

Let $K$ denote the number of ordered fit levels; throughout, $K = 3$ with
levels $\{\text{small}, \text{fit}, \text{large}\}$ encoded as
$y \in \{0, 1, 2\}$ in that order. A feature vector $\xvec \in \mathbb{R}^d$
describes one (user, item, requested size) review. The logistic function is
$\sigma(z) = (1 + e^{-z})^{-1}$. Weights are written $\wvec \in \mathbb{R}^d$
and ordered thresholds $\theta_1 < \theta_2 < \dots < \theta_{K-1}$.

\section{Introduction}
\label{sec:introduction}

Most tabular recommender portfolios that touch an explicit, multi-level
outcome treat it as an unordered multinomial classification problem, or
collapse it into a binary "good/bad" signal and reach for a ranking metric.
Clothing-fit prediction gives a clean counter-example to that default. A
garment reported as \emph{small}, \emph{fit}, or \emph{large} is not three
unrelated categories: the levels are strictly ordered, and the cost of
confusing \emph{small} for \emph{large} is qualitatively worse than
confusing \emph{small} for \emph{fit}. A model that ignores this ordering
discards information the problem handed it for free.

\subsection{Two contributions, stated precisely}

This report accompanies a working repository, not a novel method. Its two
contributions are methodological:

\begin{enumerate}
  \item \textbf{Ordinal structure, not multinomial.} We fit a
        proportional-odds cumulative link model (Section~\ref{sec:methodology})
        instead of treating the three fit levels as unordered classes, and we
        show this pays off on ordinal-aware metrics
        (Section~\ref{sec:results}) even though both models see identical
        features.
  \item \textbf{Calibrated uncertainty, not a bare label.} Wrapping the
        fitted model in split conformal prediction turns a point guess into
        a \emph{set} of plausible outcomes with a finite-sample marginal
        coverage guarantee (Section~\ref{sec:methodology}). For a sizing
        recommendation, "this could run fit or large" is a more honest and
        more actionable signal than a single confident label that is
        sometimes simply wrong.
\end{enumerate}

Both pieces are classical machine learning: the cumulative link model is fit
by maximizing a smooth, explicit likelihood with
\texttt{scipy.optimize.minimize}, and the conformal calibration is closed-form
arithmetic over a held-out calibration set. No neural network, embedding
learned by gradient descent, or deep framework appears anywhere in the
pipeline.

\subsection{Evaluation honesty}

Users are split at the group level: every review by a given user lands in
exactly one of train, validation, or test
(Section~\ref{sec:data}). This prevents the model from memorizing one user's
exact body-measurement combination and matching it back to that user's
historical fit outcomes -- a leakage mode specific to supervised tabular
problems with repeated subjects, distinct from the temporal point-in-time
leakage that recommender systems usually worry about. Both leakage modes are
present in this dataset and are tested independently
(Section~\ref{sec:feature-engineering}).

Hyperparameters for every model, baselines included, are selected on the
validation partition; the test partition is read once, to produce
Table~\ref{tab:results}.

\subsection{Repository}

The implementation and the scripts that generate every table in this
document are available in the accompanying repository. The results table
(Table~\ref{tab:results}) is emitted by
\texttt{scripts/export\_results\_table.py} and the feature catalogue
(Table~\ref{tab:features}) by \texttt{scripts/export\_feature\_tables.py}, so
the document cannot drift away from the code.

\section{Data}
\label{sec:data}

\subsection{Source}

The target dataset is the Clothing Fit Dataset for Size Recommendation
(Kaggle: \texttt{rmisra/clothing-fit-dataset-for-size-recommendation})~\cite{misra2018decomposing},
covering garment reviews from online rental and retail marketplaces. Each
review carries static shopper body measurements (height, weight, bust size,
body type, age), an item identifier and category, the purchased garment
size, and the ordinal outcome \texttt{fit} $\in \{\text{small}, \text{fit},
\text{large}\}$.

\subsection{Synthetic fixture}

\texttt{scripts/download\_data.py} defaults to a synthetic generator (no
Kaggle account required, so the full pipeline -- including CI -- runs from a
clean checkout) that reproduces the dataset's structural properties rather
than its exact values:

\begin{itemize}
  \item \textbf{Clustered users and items.} Each user has a latent "true
        size"; each item has a nominal size and a fit sensitivity. Users
        self-select garments close to their own size (as real shoppers do),
        with enough spread that mismatches still occur regularly.
  \item \textbf{Skewed fit distribution.} The probability of \emph{small}
        or \emph{large} rises smoothly with the gap between a user's true
        size and the purchased item's nominal size, scaled by the item's
        sensitivity. The resulting label distribution is skewed toward
        \emph{fit} ($56.0\%$ fit, $23.3\%$ small, $20.7\%$ large in the
        fixture used for Table~\ref{tab:results}), matching the real
        dataset's class imbalance.
  \item \textbf{True chronology.} Every review carries a \texttt{review\_date},
        so the point-in-time machinery (Section~\ref{sec:feature-engineering})
        is exercised against real ordering, not an append-order proxy.
\end{itemize}

A \texttt{--kaggle} flag downloads the real dataset via the Kaggle API when
\texttt{KAGGLE\_USERNAME}/\texttt{KAGGLE\_KEY} are set in the environment;
credentials are never read from or written to the repository.

\subsection{Grouped user split}

Train, validation and test partitions are formed by randomly assigning
\emph{users} (not rows) to a partition, so every review by a given user lands
in exactly one partition
(\texttt{recommender.data.splitting.grouped\_user\_split}). This is distinct
from -- and in addition to -- the point-in-time cut-off described in
Section~\ref{sec:feature-engineering}: it specifically prevents a model from
learning "this exact combination of height/weight/body-type always runs
small" from a user's train-partition reviews and then trivially matching it
against the same user's test-partition reviews. \texttt{GroupSplit} raises if
construction ever produces an overlapping user across partitions, and this is
asserted directly in \texttt{tests/test\_features\_leakage.py}.

\section{Feature engineering}
\label{sec:feature-engineering}

Table~\ref{tab:features} lists every declared feature, generated directly
from \texttt{conf/features.yaml} -- the single source of truth consumed by
both the training job and the serving path
(Section~\ref{sec:methodology-serving}).

% Generated by scripts/export_feature_tables.py -- do not edit by hand.
\begin{table}[t]
  \centering
  \caption{Engineered features. Static attributes need no cut-off;
  historical aggregates are computed strictly before the row being
  predicted (point-in-time).}
  \label{tab:features}
  \begin{tabular}{llll}
    \toprule
    Feature & Family & Type & Point-in-time \\
    \midrule
    \texttt{user\_height\_cm} & user & static & -- \\
    \texttt{user\_weight\_kg} & user & static & -- \\
    \texttt{user\_bmi} & user & static & -- \\
    \texttt{user\_bust\_band} & user & static & -- \\
    \texttt{user\_bust\_cup\_ordinal} & user & static & -- \\
    \texttt{user\_body\_type} & user & static & -- \\
    \texttt{user\_age} & user & static & -- \\
    \texttt{item\_category} & item & static & -- \\
    \texttt{item\_nominal\_size} & item & static & -- \\
    \texttt{user\_prior\_review\_count} & interaction & historical aggregate & \checkmark \\
    \texttt{user\_prior\_mismatch\_rate} & interaction & historical aggregate & \checkmark \\
    \texttt{user\_prior\_small\_rate} & interaction & historical aggregate & \checkmark \\
    \texttt{item\_prior\_mismatch\_rate} & interaction & historical aggregate & \checkmark \\
    \texttt{context\_rented\_for} & context & static & -- \\
    \bottomrule
  \end{tabular}
\end{table}


\subsection{Static versus historical aggregate: the distinction that matters}

This dataset makes a distinction that is easy to blur and important to get
right.

\textbf{Static attributes} -- height, weight, BMI, bust band and cup, body
type, age, item category -- are known the moment a user fills in a profile
or an item enters the catalogue. They carry \emph{no} point-in-time
cut-off because there is no "future" version of a person's height to leak:
these values simply exist before any review is written, train or test
alike. \texttt{recommender.features.user\_features} and
\texttt{item\_features.item\_category} read them directly off the row being
featurized.

\textbf{Historical aggregates} -- a user's or an item's past mismatch rate,
and a user's prior review count -- are a fundamentally different kind of
feature: they summarize \emph{outcomes of other reviews}, and the review
being predicted is one such outcome. If a user's "prior mismatch rate" were
computed over \emph{all} of that user's reviews (including the one being
predicted, or ones that happened later), the feature would partially encode
the label it is meant to help predict. This is exactly the leak this project
is built to avoid.

\subsection{Point-in-time construction}

Each review carries two chronological positions, assigned once when the
dataset is indexed (\texttt{recommender.data.loader.build\_indexed\_dataset},
\texttt{recommender.features.pipeline.assign\_global\_order}):

\begin{itemize}
  \item \texttt{review\_order} -- the review's position within its own
        user's history (0, 1, 2, \ldots), used by the two user-level
        aggregates.
  \item \texttt{global\_order} -- the review's position within the entire
        dataset's chronology (by \texttt{review\_date}), used by the
        item-level aggregate, since an item can be reviewed by many
        different users across every partition.
\end{itemize}

For a review at position $k$ (per the relevant order column), the historical
aggregate is an \emph{expanding} statistic over strictly prior positions
only:
\begin{equation}
  \widehat{r}(k) = \frac{1}{k} \sum_{j < k} \mathbb{1}[\text{outcome}_j
  \text{ was a mismatch}], \qquad \widehat{r}(0) \text{ undefined (no prior
  history)}.
\end{equation}
The strict inequality is the entire contract. A review at a given position
can see every review strictly before it, never the review at its own
position, and never one after it. This is enforced by computing the
statistic over a frame sorted by the order column and using an explicit
less-than comparison, not a window function that could silently include the
current row (\texttt{interaction\_features.\_expanding\_prior\_mean}).

Rows with no qualifying prior history (a user's or item's very first
appearance) get \texttt{NaN} from the raw computation; the feature pipeline
then applies the fill policy declared in \texttt{conf/features.yaml} -- the
global mismatch rate for the two rate features, zero for the count and the
small-rate feature, since a user's first review has, by definition, no
observed directional bias yet.

\subsection{Leakage tests}

\texttt{tests/test\_features\_leakage.py} operationalizes the point-in-time
contract directly: it injects a fabricated future review (an extreme, label-
flipping outcome) for an existing user or item, strictly after the real
cut-off, and asserts that every feature value for rows at or before that
cut-off is unchanged. A second test asserts the point-in-time aggregate
differs from a naive whole-history average whenever a user has more than one
review -- if they ever coincided, every row would be seeing its own future.

\subsection{Context features}

A shopper-stated occasion (\texttt{context\_rented\_for}, e.g.\ "wedding",
"work") is supplied at request time, like a search filter. It is information
about the current transaction, not a historical aggregate, so it needs no
cut-off either -- it exists in exactly the same sense static attributes do.

\subsection{Training/serving parity}

\texttt{FeatureStore.transform} (batch) and \texttt{FeatureStore.row\_vector}
(single serving-time row) are the only two entry points into the feature
pipeline, and \texttt{tests/test\_features\_parity.py} asserts they produce
numerically identical output for the same logical row -- the failure mode
this guards against is an offline batch computing one feature vector while
an online request computes a different one for what should be the same
input.

\section{Methodology}
\label{sec:methodology}

\subsection{Cumulative link (proportional-odds) model}

For an ordinal outcome with $K$ levels and feature vector $\xvec$, the
cumulative-link model~\cite{mccullagh1980,agresti2010} with a logistic link
assumes

\begin{equation}
  P(y \le k \mid \xvec) = \sigma(\theta_k - \wvec^\top \xvec),
  \qquad k = 1, \dots, K - 1,
  \label{eq:cumulative-link}
\end{equation}

with ordered thresholds $\theta_1 < \theta_2 < \dots < \theta_{K-1}$ and a
single weight vector $\wvec$ shared across every threshold -- the
"proportional odds" / "parallel lines" assumption: a one-unit increase in any
feature shifts every cumulative log-odds by the same amount
$-\wvec^\top \xvec$, which is what makes the model genuinely ordinal rather
than $K-1$ unrelated logistic regressions. By convention
$P(y \le 0 \mid \xvec) = 0$ and $P(y \le K - 1 \mid \xvec) = 1$, so the
per-class probability is the first difference of the cumulative
distribution,
\begin{equation}
  P(y = c \mid \xvec) = P(y \le c \mid \xvec) - P(y \le c - 1 \mid \xvec),
  \qquad c = 0, \dots, K - 1.
  \label{eq:class-prob}
\end{equation}

\textbf{Unconstrained parametrization.} An optimizer searching over
$(\theta_1, \dots, \theta_{K-1})$ directly could produce a non-monotonic,
invalid threshold vector. We instead optimize over an unconstrained vector
$\boldsymbol{\delta} \in \mathbb{R}^{K-1}$ and set
\begin{equation}
  \theta_1 = \delta_1, \qquad
  \theta_k = \theta_{k-1} + \operatorname{softplus}(\delta_k)
  \;\; (k = 2, \dots, K-1),
\end{equation}
with $\operatorname{softplus}(z) = \log(1 + e^z) > 0$. Because every gap
after the first is a strictly positive softplus output, $\theta_1 < \theta_2
< \dots < \theta_{K-1}$ holds for \emph{any} real-valued $\boldsymbol{\delta}$,
so an unconstrained quasi-Newton optimizer can never propose an invalid
threshold ordering. This is implemented in
\texttt{recommender.models.ordinal.\_thresholds\_from\_raw}.

\subsection{Negative log-likelihood objective}

Given training data $\Dset = \{(\xvec_i, y_i)\}_{i=1}^{n}$, parameters
$\wvec$ and $\boldsymbol{\theta}$ are fit by minimizing the $\ell_2$-penalized
negative log-likelihood
\begin{equation}
  \mathcal{L}(\wvec, \boldsymbol{\theta})
  = -\sum_{i=1}^{n} \log P(y = y_i \mid \xvec_i; \wvec, \boldsymbol{\theta})
  \; + \; \frac{\lambda}{2} \lVert \wvec \rVert_2^2 ,
  \label{eq:nll}
\end{equation}
with $P(y = y_i \mid \cdot)$ given by Equation~\ref{eq:class-prob}
substituted into Equation~\ref{eq:cumulative-link}. The thresholds are not
penalized, since shrinking a cut-point toward zero has no natural
interpretation the way shrinking a feature weight does. Equation~\ref{eq:nll}
is minimized directly with \texttt{scipy.optimize.minimize} (L-BFGS-B) over
the unconstrained parameter vector
$(\delta_1, \dots, \delta_{K-1}, \wvec)$ -- classical convex-leaning
optimization, with no gradient-descent-trained embedding or neural layer
anywhere in the model
(\texttt{recommender.models.ordinal.OrdinalLogisticRegression}).

\subsection{Baselines}

\paragraph{Majority class.} Predicts the single most frequent training
label for every row, ignoring features entirely. The floor any model using
features must clear.

\paragraph{Multinomial logistic regression.} The same engineered features,
fed into an \emph{unordered} softmax classifier,
\begin{equation}
  P(y = c \mid \xvec) = \frac{\exp(\wvec_c^\top \xvec)}
  {\sum_{c'=0}^{K-1} \exp(\wvec_{c'}^\top \xvec)},
\end{equation}
fit by the same $\ell_2$-penalized NLL minimization via
\texttt{scipy.optimize}. Because it assigns an independent weight vector
$\wvec_c$ to every class, it cannot exploit the ordering the way the
proportional-odds model's shared $\wvec$ does -- isolating the ordinal
structure's effect, rather than any difference in optimizer or feature set,
as the variable under test (Section~\ref{sec:results}).

\subsection{Split conformal prediction}

Any fitted classifier that outputs $\hat{P}(y \mid \xvec)$, ordinal or not,
can be wrapped in split conformal prediction to produce a prediction
\emph{set} with a finite-sample marginal coverage guarantee -- implemented
from scratch in NumPy in \texttt{recommender.models.conformal}.

\paragraph{Nonconformity score.} On a calibration set
$\Dset_{\text{cal}} = \{(\xvec_i, y_i)\}_{i=1}^{m}$, disjoint from both the
training and test partitions (here: the validation partition), the
nonconformity score for example $i$ is
\begin{equation}
  \alpha_i = 1 - \hat{P}(y_i \mid \xvec_i),
  \label{eq:nonconformity}
\end{equation}
large exactly when the fitted model placed little probability mass on the
true label.

\paragraph{Calibration procedure.} For a target miscoverage rate
$\alpha \in (0, 1)$, let
\begin{equation}
  \hat{q} = \text{the } \left\lceil (m+1)(1-\alpha) \right\rceil / m
  \text{-th empirical quantile of } \{\alpha_1, \dots, \alpha_m\}.
  \label{eq:qhat}
\end{equation}
For a new point $\xvec$, the conformal prediction set is
\begin{equation}
  C(\xvec) = \left\{ y \in \{0, \dots, K-1\} :
  1 - \hat{P}(y \mid \xvec) \le \hat{q} \right\}.
  \label{eq:prediction-set}
\end{equation}

\paragraph{Coverage guarantee.} Under the sole assumption that the
calibration example and the test example are exchangeable (a strictly
weaker assumption than any distributional or parametric one),
\begin{equation}
  P\big(y \in C(\xvec)\big) \; \ge \; 1 - \alpha,
  \label{eq:coverage-guarantee}
\end{equation}
marginally over the randomness in both the calibration set and the new test
point~\cite{vovk2005,lei2018,angelopoulos2023}. This guarantee is model-agnostic: it holds
regardless of whether $\hat{P}$ is well-calibrated, over-confident, or
under-confident, which is precisely why conformal prediction is the right
tool to pair with a point-prediction model whose own calibration is not
otherwise guaranteed.

\subsection{Why this problem is calibration, not ranking}
\label{sec:why-not-ranking}

NDCG, Recall@K and MAP score a \emph{ranked list of items} against a known
set of relevant ones per user, under an implicit assumption that the
consumer will see the list top-to-bottom. Nothing in this problem has that
shape: every example is one independent (user, item, requested size) review
with a single explicit ordinal label, not a list to be cut off at $K$. There
is no notion of "relevant items for this user" to rank against, and no
top-$K$ truncation that would make a ranking metric meaningful. We therefore
report ordinal-aware classification and calibration metrics
(Section~\ref{sec:evaluation}) instead of adapting a ranking metric that does
not fit the problem's shape.

\subsection{Serving}
\label{sec:methodology-serving}

The fitted model, its scaler, and its calibrated conformal wrapper are
applied identically whether scoring a batch (training/evaluation) or a
single incoming request
(\texttt{recommender.serving.predictions.build\_prediction\_payloads}).
Rather than running inference inside the request path, predictions for every
anticipated \texttt{(user, item, size)} triple are precomputed and stored in
DynamoDB keyed by a composite request id; the inference Lambda
(\texttt{recommender.serving.handler.lambda\_handler}) is then a pure key
lookup, bounding p99 latency by DynamoDB rather than by feature
reconstruction cost.

\section{Evaluation protocol}
\label{sec:evaluation}

\subsection{Partitioning}

Users are randomly assigned to train ($70\%$), validation ($15\%$) and test
($15\%$), and every review by a given user lands in exactly one partition
(\texttt{grouped\_user\_split}, Section~\ref{sec:data}). The split is
reconstructed and checked for user overlap on every call
(\texttt{GroupSplit.\_\_post\_init\_\_}), and
\texttt{tests/test\_features\_leakage.py} asserts that attempting to build a
split from row-sliced (rather than user-grouped) partitions raises.

\subsection{Metrics}

\paragraph{Mean absolute ordinal error.} With $y_i \in \{0, 1, 2\}$ the true
level and $\hat{y}_i$ the predicted level,
\begin{equation}
  \text{MAE}_{\text{ord}} = \frac{1}{n} \sum_{i=1}^{n} \lvert y_i - \hat{y}_i \rvert.
\end{equation}
Unlike plain accuracy, this distinguishes a \emph{small}-predicted-as-\emph{large}
error (distance 2) from a \emph{small}-predicted-as-\emph{fit} error
(distance 1) -- the entire reason to treat the problem as ordinal.

\paragraph{Accuracy within one level.} The fraction of predictions at most
one ordinal step from the truth, $\frac{1}{n}\sum_i
\mathbb{1}[\lvert y_i - \hat{y}_i\rvert \le 1]$ -- a tolerant metric that
still distinguishes a model that is occasionally off by one level from one
that confuses \emph{small} and \emph{large} outright.

\paragraph{Macro F1.} Unweighted mean of per-class F1, reported because
accuracy alone is a weak signal under the dataset's class skew (the majority
baseline attains high accuracy and within-one-level scores by predicting
"fit" for everyone, but a macro F1 near the single-class floor).

\paragraph{Conformal coverage and average set size.} Empirical fraction of
test examples whose true label is included in the conformal set
(Equation~\ref{eq:prediction-set}), and the mean set cardinality. Coverage
should sit near the target $1-\alpha$; set size is the cost of that
guarantee -- a trivial predictor could "achieve" perfect coverage by
returning every class every time (as the majority baseline does in
Table~\ref{tab:conformal}), which is why set size must always be read
alongside coverage, not instead of it.

\paragraph{What we do not report, and why.} NDCG@K, Recall@K and MAP are
absent from this report by design, not omission; see
Section~\ref{sec:why-not-ranking} for the argument.

\subsection{Model selection and test discipline}

Every model's hyperparameters ($\ell_2$ strength for both logistic models)
were selected by inspecting validation-partition metrics; the test partition
is touched only once, by \texttt{scripts/export\_results\_table.py}, to
produce Table~\ref{tab:results}. The conformal calibration set is the
validation partition, disjoint from both the training data the model's
\emph{parameters} were fit on and the test data its \emph{coverage} is
measured on -- using the same partition for fitting the model and
calibrating the conformal wrapper would violate the exchangeability
assumption the coverage guarantee (Equation~\ref{eq:coverage-guarantee})
depends on.

\subsection{Training/serving parity}

\texttt{tests/test\_serving.py} asserts every precomputed payload's predicted
class is a member of its own conformal set (the fitted model's own top
choice can never be excluded by its own calibrated threshold), that
class probabilities sum to one, that request ids are unique, and that the
full payload is JSON-serializable -- the same check the reference
BPR/SVD portfolio projects use to catch NumPy-scalar serialization failures
before they reach a Lambda deployment.

\section{Results}
\label{sec:results}

% Generated by scripts/export_results_table.py -- do not edit by hand.
\begin{table}[t]
  \centering
  \caption{Point-prediction quality on the held-out test partition. Hyperparameters for every model were selected on the validation partition; the test partition was read once.}
  \label{tab:results}
  \begin{tabular}{lrrrr}
    \toprule
    Model & Accuracy & MAE (ordinal) & Acc. within 1 & Macro F1 \\
    \midrule
    OrdinalLR & 0.6509 & \textbf{0.3535} & 0.9956 & 0.5635 \\
    Majority & 0.5603 & 0.4397 & 1.0000 & 0.2394 \\
    MultinomialLR & 0.4258 & 0.6298 & 0.9444 & 0.4283 \\
    \bottomrule
  \end{tabular}
\end{table}


Table~\ref{tab:results} reports point-prediction quality on the held-out test
partition. The ordinal model attains the lowest mean absolute ordinal error
by a wide margin over both baselines, confirming that the proportional-odds
structure is doing real work: the multinomial baseline sees identical
features but, lacking the shared-weight ordinal constraint, distributes its
probability mass less efficiently across adjacent levels and pays for it
on every ordinal-aware metric. The majority baseline's accuracy-within-one
score of $1.0000$ is a direct consequence of always predicting "fit" (the
dataset's modal class) under a tolerance of one level either side --
illustrating exactly why within-one accuracy must be read alongside mean
absolute error and macro F1, not instead of them: its macro F1 of
$0.2394$ is the lowest of the three models, reflecting that it is correct
only on the majority class.

% Generated by scripts/export_results_table.py -- do not edit by hand.
\begin{table}[t]
  \centering
  \caption{Split conformal calibration at target coverage $1-\alpha=0.9$, calibrated on the validation partition and evaluated on test.}
  \label{tab:conformal}
  \begin{tabular}{lrrr}
    \toprule
    Model & Target cov. & Empirical cov. & Avg. set size \\
    \midrule
    Majority & 0.9000 & 1.0000 & 3.0000 \\
    MultinomialLR & 0.9000 & 0.8922 & 2.0784 \\
    OrdinalLR & 0.9000 & 0.8855 & 1.7154 \\
    \bottomrule
  \end{tabular}
\end{table}


Table~\ref{tab:conformal} reports the split conformal calibration at a
$90\%$ target coverage, calibrated on the validation partition and
evaluated on test. The ordinal model's conformal sets attain empirical
coverage close to the $90\%$ target while remaining the most informative:
an average set size of $1.72$ out of $3$ possible classes, versus $2.08$
for the multinomial baseline and a trivial $3.0$ (every class, always) for
the majority baseline. The majority baseline's set size directly reflects
that its point probabilities are too flat (identical for every input) to
ever confidently exclude a class -- it "achieves" coverage only by refusing
to narrow the answer at all, the degenerate case the average-set-size metric
exists to catch.

\subsection{Reading the headline number}

The practical takeaway is not "$65\%$ accuracy," which undersells what the
pipeline actually produces. It is: for a held-out shopper the system never
saw during training, the calibrated output is typically a set like
$\{\text{fit}\}$ or $\{\text{fit}, \text{large}\}$ with a guarantee that the
true outcome falls inside that set at least $90\%$ of the time -- a result a
single point prediction cannot make and cannot be checked against.

\section{Conclusion}
\label{sec:conclusion}

We treated clothing-fit prediction as what it structurally is -- an ordinal
classification problem over three ordered levels -- rather than defaulting
to an unordered multinomial classifier or collapsing it into a binary
signal, and paired the fitted model with split conformal prediction so its
output is a calibrated set, not a bare label. Both pieces are classical
machine learning: a cumulative link model fit by maximizing an explicit
likelihood with \texttt{scipy.optimize}, and a closed-form calibration
procedure over a held-out partition, implemented from scratch in NumPy. No
neural network, embedding layer, or gradient-descent training loop appears
anywhere in the pipeline.

The ordinal structure measurably helped: against an unordered multinomial
logistic regression trained on identical features, the proportional-odds
model nearly halved the mean absolute ordinal error ($0.3535$ versus
$0.6298$) and produced visibly tighter conformal sets ($1.72$ versus $2.08$
classes on average) at the same target coverage. Against the naive
majority-class floor, it traded a small amount of macro F1 generosity to
the majority class for a model that actually uses the shopper's and
garment's features, reflected in more than double the majority baseline's
macro F1.

\subsection{Limitations}

\textbf{Synthetic fixture.} All numbers in this report are computed on the
synthetic fixture, not the real Kaggle dataset; \texttt{scripts/download\_data.py
--kaggle} provides the path to real data when credentials are available, and
the feature pipeline, model, and tests are schema-driven rather than tuned to
the fixture's specific distribution.

\textbf{Proportional-odds assumption.} The shared-weight cumulative link
model assumes every feature's effect on the cumulative log-odds is constant
across thresholds. If, for example, body type affected the \emph{small}
threshold very differently from the \emph{large} threshold, a
non-proportional (threshold-varying) extension would fit better; we did not
find evidence requiring that extra complexity on this fixture, but a
partial-proportional-odds variant is the natural next step on real data.

\textbf{Conformal set size on a 3-class problem.} With only three classes,
an "informative" conformal set can only take sizes $\{1, 2, 3\}$; the method
scales more informatively to problems with more ordinal levels (e.g. a 5-star
or 10-point sizing scale), where a tight conformal set communicates
substantially more than it can here.

\subsection{What a reader should take away}

The point this project makes is not that an ordinal model is marginally more
accurate. It is that treating an explicitly ordered outcome as ordered, and
reporting a model's own uncertainty honestly instead of suppressing it behind
a single confident label, are both small implementation choices with an
outsized effect on whether a recommendation can be trusted.


\bibliographystyle{plain}
\bibliography{references}

\end{document}
